% ============================================
% MEM-04 (ZH)
% 中文翻译 — SPEC-935-R356
% ============================================
\section*{\textbf{MEM-04} --- 乱葬坑}
\markboth{MEM-04 --- 乱葬坑}{}

\noindent \textit{相遇}

\noindent \textbf{1915 —— 从塞纳多尔·蓬佩乌到福塔莱萨的路上}

\vspace{0.5em}\noindent\rule{\textwidth}{0.5pt}\vspace{0.5em}

\noindent 我父亲是在早晨死的。太阳升起时，他已经僵硬了，眼睛睁着，望向没有云的天空。母亲用手合上他的眼睛，说：\textit{“走吧，若阿金。你爹已经走了。”}

\noindent 那年我八岁。

\noindent 我们已经走了四天。以前是六个人。现在只剩三个：我、母亲，和安塞尔莫老爹——他不是亲戚，但和我们一起走，因为独自一人谁也活不下去。安塞尔莫老爹说他见过1877年的湖干涸。说那才是真正的旱。我看着他，看出连他自己也不信自己说的话。他的眼睛和我父亲的一模一样。

\noindent 第五天，我们遇到了那队人。

\noindent 是士兵。瘦弱、肮脏，穿着打补丁的制服。一个背着步枪。另一个拿着鞭子。

\noindent \textit{“去哪儿？”拿鞭子的那个问道。}

\noindent 母亲说我们要去海边。说她听说那里有食物。

\noindent 士兵笑了。\textit{“食物？海边挤满了像你这样的逃荒者。大家都在海滩上等死。”}

\noindent 母亲没有回答。

\noindent \textit{“跟我们走，”士兵说，“福塔莱萨附近的阿拉加迪索（Alagadiço）有一个收容营。那里有吃的。”}

\noindent 她犹豫了。我感觉到她的手握紧了我的手。

\noindent \textit{“这是命令，”士兵说，鞭子在空气中啪地一响。没有打到任何人。但那声音留下了。仿佛在我们身上烙下了印。}

\vspace{0.5em}\noindent\rule{\textwidth}{0.5pt}\vspace{0.5em}

\noindent 政府牲畜圈是一圈围着空地的铁丝网。没有树。没有荫。有人。很多人。比我一生中见过的所有人都多。

\noindent 他们躺在地上，彼此靠着。有些人动着嘴却发不出声音。有些人一动不动。

\noindent \textit{“莫兰布多斯（破衣人）”，有人说。}

\noindent 我不知道那是什么意思。后来我知道了。意思是破布。人形的破布。不再是人的人，因为干旱带走了他们体内的一切。

\noindent 我们留了下来。母亲在铁丝网边找到一块地方。安塞尔莫老爹坐下，就再也没有站起来。第二天，他说：\textit{“我要走了。”}然后他就走了。

\noindent 母亲开始咳嗽。

\vspace{0.5em}\noindent\rule{\textwidth}{0.5pt}\vspace{0.5em}

\noindent 乱葬坑在牲畜圈后面，大约两百米远。每天晚上士兵们把尸体搬过去。白天能看见秃鹫。它们在天上绕圈，转啊转，仿佛天空是一只巨大的眼睛。

\noindent 在一个母亲咳个不停的夜晚，我走到了乱葬坑边。

\noindent 我不是出于好奇去的。
我去是因为听见了一个东西。

\noindent 是一个声音。

\noindent 不是人的声音——但确实是声音。它在唱歌。不是有词的歌。是一段旋律。低沉。沙哑。一段像风穿过空地的旋律，起起伏伏。

\noindent 它唱着：

\noindent \textit{“饿的人坐下。饿的人等待。饿的人有什么吃什么。”}

\noindent 我在坑边停下。

\noindent 下面，在尸体之间，有什么东西在动。

\noindent 我本该跑掉。本该回到牲畜圈。本该把脸埋进母亲的裙子里，假装什么都没听见。

\noindent 但我只是个孩子。
而饿着的孩子不会逃避食物。
他们会去查看。

\vspace{0.5em}\noindent\rule{\textwidth}{0.5pt}\vspace{0.5em}

\noindent 坑不深。大概三米。尸体像面粉袋一样堆叠着。有些是赤裸的。有些双手被绑着。他们的眼睛睁着。所有的眼睛都睁着。

\noindent 在尸堆中间，有什么东西小心地移动着。

\noindent 那不是挣扎着要爬出来的动作。
那是正在\textit{整理}周围尸体的动作。

\noindent 一只手臂伸了出来。不是人的手臂——太细了，太长了，关节长在不该长的地方。手上的手指太多，每根手指都有一枚黑色的指甲，像湿玻璃一样发亮。

\noindent 那只手抓住尸堆中一条松脱的手臂，把它放到一边。很小心。像人在用餐前摆桌子。

\noindent 另一只手臂出现了。然后是头。

\noindent 那头几乎像人的。几乎。形状是骷髅，但皮肤还裹着骨头——一层薄薄的、灰色的皮，像湿纸一样绷在面骨上。

\noindent 鼻子是一个洞。
嘴是一条线。
眼睛是两道缝。

\noindent 而在缝隙里，有一抹黄色的光。

\noindent 它看着我。

\noindent 没有动。什么也没做。只是看着我，歪着头，像一条狗在听陌生的声音。

\noindent 我没有跑。

\noindent 它开始爬上来。

\noindent 不像人那样爬。像蜘蛛。手臂和腿协调地移动，在尸体上寻找支撑，不紧不慢地在它们之间滑行。骨头在它的重量下咔咔作响。肉塌陷下去。

\noindent 它爬出了坑。

\noindent 站在了我面前。

\noindent 它巨大无比——不是高，而是\textit{宽}。身体是一架朝外翻的肋骨笼子，仿佛骨架长反了。肋骨之间，有什么东西在搏动。一个巨大的、黑暗的器官，布满像蚯蚓一样发亮的血管。

\noindent 它的气味飘到我面前。

\noindent 不是死亡的气味。是\textit{食物}的气味。烤肉。浓汤。旱灾之前辛哈太太厨房里飘出的气味。我已经好几个月没有闻到的气味。

\noindent 我的胃疼了起来。

\noindent \textit{“你饿吗？”}

\noindent 它的声音很柔。几乎甜美。和身体毫不相称。像是祖母在问孙子吃过饭没有。

\noindent 我没有回答。

\noindent \textit{“你饿。”它把头歪向另一边。“我感觉得到。你的饿很大。大得跟我的饿一样。”}

\noindent \textit{“您……”我的声音很弱。“您是什么？”}

\noindent 它想了想。肋骨之间的器官搏动得更厉害了。

\noindent \textit{“我是剩下的东西，”它说，“我是很多死亡之处生长出来的东西。我是大地咽不下去的东西。”}

\noindent \textit{“您吃死人吗？”}

\noindent \textit{“总得有人吃。”}

\noindent 它说这话时那么自然。像在说“太阳每天都会升起”。

\noindent \textit{“活人不愿意。活人埋葬。活人哭泣。然后离开。死人留下来。总得有人利用。”}

\noindent 它向我伸出手。那只大手，手指修长，指甲漆黑。

\noindent \textit{“你要吗？”}

\noindent \textit{“要什么？”}

\noindent \textit{“吃。”}

\noindent 我看着它的手。看着坑。看着那些尸体。

\noindent \textit{“不，”我说，“我不想吃他们。”}

\noindent 它笑了。笑声像干枯的叶子在翻动。

\noindent \textit{“我给你的不是他们，孩子。”}

\noindent \textit{“那是什么？”}

\noindent 它弯下腰。脸凑到和我一样高。食物的气味更浓了。我的胃又疼起来。

\noindent \textit{“我替你吃，”它说，“我每天都吃。我吃死去的人。我吃没有人要的人。而你，再也不会感到饥饿。”}

\noindent 它等着。

\noindent \textit{“只是……”}

\noindent \textit{“只是什么？”}

\noindent 它笑了。嘴张得比应该的大。牙齿排成了太多排。多得装不进一张人的嘴。

\noindent \textit{“饥饿不会离开。它留下来。只不过它不再是你的了。它是我的。而我要住进你体内，去感受它。”}

\noindent 它用一根手指碰了碰我的胸口。

\noindent \textit{“你接受吗？”}

\noindent 我本该说不。

\noindent 我本该朝它的手吐唾沫。本该跑掉。本该尖叫。

\noindent 但我只有八岁。
而且我那么饿。
饿得那个“不”字怎么也出不来。

\noindent \textit{“好，”我说。}

\noindent 它张开了嘴。

\noindent 不是亲吻。不是撕咬。是一口气。一股温热的气从它嘴里出来，进入我的嘴。滑下喉咙。灌满胸腔。淌进胃里。

\noindent 它在我体内了。

\noindent \textit{“成了，”它说，“现在饥饿是我的了。”}

\noindent 它站起来。望着天空。月亮亮着。

\noindent \textit{“去吧，”它说，“去活着。以后我们还会再见。”}

\noindent 它回到坑里。在尸体之间下去。躺下。闭上眼睛。

\noindent 我站在坑边，看着。

\noindent 我不再感到饥饿。
但感到一种沉重。
一种沉在胃底的重量。
一个会呼吸的重量。\NE{与乱葬坑订约的段落，是整本日记中唯一一处若阿金用左手写下的文字。检验过手稿的专家注意到：在“好，我说”那一行，笔迹突然改变，笔尖对纸面的压力不同，字母的走向向左偏斜——仿佛另一只手在引导着他的手。编辑没有干预。保留原始笔迹的选择，是一种伦理和文献上的选择。}

\vspace{0.5em}\noindent\rule{\textwidth}{0.5pt}\vspace{0.5em}

\noindent 那一夜，在一个死人的乱葬坑边，在一片集中营中央，一个八岁的男孩与一个怪物做了交易。

\noindent 怪物吃掉了他饥饿。

\noindent 但留下了那个洞。

\noindent 洞留了下来。

\noindent 此后进入洞中的一切——岁月、血、人、记忆——永远也填不满。

\noindent 因为那个洞不是我的。

\noindent 是它的。

\noindent 从那夜起，它就住在我体内。

\noindent 它吃。我感觉。

\noindent 它睡。我醒来。

\noindent 它等待。我活着。

\noindent 它也在等你。

\vspace{0.5em}
\noindent\textit{\textbf{↪ 路线:} \rota{MEM-05} \rotasep \rota{DOC-01} \rotasep \rota{LUC-01}}
